% 本文件由 LaTeX 排版 Agent 根据有效 translation.json 自动生成。
% author=Apocrypha; work_id=0818; title=斐理伯行传

\chapter{斐理伯行传}

% structure: p0001 source=h1 target=omitted reason=page-title-already-rendered-by-parent

% structure: p0002 source=h2 target=section reason=relative-heading-rank; base=h2

\section{论宗徒斐理伯的旅程，从第十五件事到末尾，其中包含他的殉道。}

大约在皇帝\Person{图拉真}开始统治罗马人时，克洛帕的儿子\Person{西满}——即耶路撒冷的主教——在他统治的第八年殉道之后，他是继被称为主兄弟的\Person{雅各伯}之后，那里教会的第二位主教，宗徒\Person{斐理伯}走遍\Place{里狄雅}和\Place{亚细亚}的各城各地区，向众人宣讲基督的福音。\par

他来到\Place{奥菲奥吕玛}城，即\Place{亚细亚}的\Place{希拉波利斯}，受到一位名叫\Person{斯塔希}的信徒接待。与他同行的还有主的七十位门徒之一的\Person{巴尔多禄茂}、他的姊妹\Person{玛利亚默}，以及跟随他的门徒们。于是全城的人都放下工作，跑到\Person{斯塔希}家里，听说斐理伯所行的事。许多男女聚集在\Person{斯塔希}家中，斐理伯同巴尔多禄茂一起教导他们关于耶稣的事。\par

斐理伯的姊妹\Person{玛利亚默}坐在\Person{斯塔希}家的门口，向来到的人说话，劝他们听宗徒们的话，对他们说：\Quote{我们兄弟，我在天之父的儿女们，你们是上等财富，是天上之城的实质，是天主为爱祂的人所预备的居所的喜乐。要践踏仇敌的陷阱，那蜿蜒的蛇。因为他的路是弯曲的，他是恶者的儿子，邪恶的毒在他内；他的父亲是魔鬼，死亡的创造者，他的母亲是败坏；愤怒在他眼中，毁灭在他口中，他的路是阴府。所以，要逃离那没有实质的、无形状的、在一切受造物中没有形状的那位，无论在天上或地上，无论飞鸟或走兽。因为一切都被从他的形状中夺走了；因为在地的走兽和天的飞鸟中，有对他的认识：那蛇拖着肚腹和胸膛；塔尔塔罗斯是他的居所，他在黑暗中行走，因为他无所倚靠。所以，要逃离他，免得他的毒注入你的口。倒要成为有信德、圣洁、行善、无诡诈的。你们要除去邪恶的性情，即那邪恶的欲望，藉此蛇，那邪恶的龙，邪恶的首领，为灵魂产生了毁灭和死亡的牧场，因为一切邪恶的欲望都从他而来。这是邪恶的根源，恶行的维持，灵魂的死亡：因为仇敌的欲望武装起来攻击信徒，从黑暗中出来，在黑暗中行走，试图与在光明中的人争战。因为这是贪欲的开端。因此，你们愿意来到我们这里，更因天主藉我们如同父亲到自己的儿女一样来到你们这里，愿意怜悯你们，拯救你们脱离仇敌的邪恶陷阱，要逃离仇敌的邪恶情欲，彻底将它们从你们心中驱逐，公开憎恨邪恶之父，并爱耶稣，祂是光、生命、真理，以及一切渴望祂者的救主。所以，要奔向祂，用爱抓住祂，使祂能把你们从恶者的坑中拉上来，并洁净你们，使你们无可指摘，生活在真理中，在祂圣父的面前。}\par

这一切都是斐理伯向那些像古时一样聚集来敬拜蛇和毒蛇的群众所说的，他们也为此立像敬拜。因此他们称希拉波利斯为奥菲奥吕玛。斐理伯说了这些话后，\Person{巴尔多禄茂}、\Person{玛利亚默}、他的门徒们，以及与他同在的\Person{斯塔希}，全体百姓都留心听，他们中许多人逃离仇敌，转向耶稣，加入了斐理伯和与他在一起的人。信徒们更加在基督的爱里被坚固。\par

\Person{尼卡诺拉}是总督的妻子，卧病在床，患多种疾病，尤其是眼睛。她听说了宗徒\Person{斐理伯}和他的教导，就信了主。因为她此前就已听说过祂；她呼求了祂的名，就从折磨她的痛苦中解脱了。于是起身，从侧门出去，由自己的奴隶用银轿抬着，来到\Person{斯塔希}家中——宗徒们就在那里。\par

当她来到门口时，宗徒\Person{斐理伯}的姊妹\Person{玛利亚默}看见她，在\Person{斐理伯}、\Person{巴尔多禄茂}和所有已信的人面前，用希伯来语对她说：\Quote{Alemakan, ikasame, marmare, nachaman, mastranan, achaman；意思是：父亲的女儿，你是我的主母，你被当作抵押交给了蛇；但我们的救主耶稣已来藉着我们救你，折断你的锁链，切断它们，将它们从你身上连根拔除，因为你是我姊妹，同一母亲生了我们二人。你已离弃你父亲，你已离弃那引你到你母亲居所的路，处于迷误中；你已离开那欺骗和暂时荣耀的殿，来到我们这里，逃离仇敌，因为他是死亡的居所。看，如今你的救赎主已来救赎你；基督，正义的太阳，已升起光照你。}\BibleRef{拉 4:2}\par

\Person{尼卡诺拉}站在门前听见这些话，就在众人面前鼓起勇气，喊叫说：\Quote{我是希伯来人，是希伯来人的女儿；请用我祖先的语言对我说话。因为我听了祖先的宣讲，立刻从疾病和包围我的痛苦中痊愈了。因此，我敬仰天主的良善，祂使你们被掠到这座城来，是为了祂那受尊崇的真石，好叫我们藉着你们认识祂，并因信祂而能与你们同生。}\par

尼卡诺拉这样说了之后，宗徒斐理伯连同巴尔托羅买、玛利亚姆和与他们一起的人，为她向天主祈祷说：\Quote{你使死人复活，主耶稣基督，你藉着圣洗使我们从死亡的奴役中获得自由（见：\BibleRef{罗 6:3}），也完全拯救这位妇女脱离那仇敌的错谬；在你的生命中使她存活，并在你的成全中使她成全，使她得以在她先祖的国土上自由地找到，在你美善中拥有份，主耶稣。}\par

所有人与宗徒斐理伯一起高呼阿们，看，来了那个暴君，尼卡诺拉的丈夫，像一匹未上辔头的马一样暴怒；他抓住妻子的衣裳，喊着说：\Quote{啊，尼卡诺拉，我不是把你留在了床上吗？你怎么有这么大的力气跑到这些术士这里来了？你的眼炎又怎么治好的？现在，你若不告诉我你的医生是谁，他叫什么名字，我就要用各式各样的刑罚惩罚你，决不会怜悯你。}她回答说，对他说：\Quote{暴君啊，从你身上除掉这暴虐，离弃这恶行；抛弃这只能持续片刻的生命；逃离你这无益禀性的野蛮；逃避那恶龙及其欲望；从身上扔掉那杀人之蛇的作为和毒箭；弃绝那可憎恶的偶像的祭献，那是敌人的耕作，黑暗的篱笆；为你自己营造贞洁纯净的生活，好能在圣洁中认识我的医生，得知他的名字。因此，你若想我在你身边，就要准备好活在贞洁和节制中，并敬畏真天主，这样我才能与你共度一生；只要你从偶像和它们的一切污秽中洁净自己。}\par

她那阴郁的暴君丈夫听了她这些话，就抓住她的头发，踢着她，一路拖着她，说：\Quote{你被刀剑砍掉，或者我看见你离开我跟这些外来的术士行淫，这对你倒是件好事；因为我看出你已堕入这些骗子的疯狂之中。因此，你，他们中的第一个，我要用邪恶的死亡除掉；然后，不饶恕他们，我要砍断他们的筋，用最残酷的死亡处死他们。}他转过身来，对身边的人说：\Quote{把那几个妖言惑众的术士给我带出来。}公差跑进斯塔齐的家，抓住宗徒斐理伯、巴尔托羅买和玛利亚姆，拖着他们，带到总督那里。最忠信的斯塔齐和所有信众跟随着。\par

总督看见他们，咬牙切齿地说：\Quote{折磨这些骗子，他们欺骗了许多妇女、青年和少女，说自己是敬拜天主的人，而他们却是可憎之物。}他下令拿来生牛皮条，鞭打斐理伯、巴尔托羅买和玛利亚姆；他们被牛皮条鞭打后，又下令捆住他们的脚，拖着他们穿过城里的街道，直拖到他们神庙的门口。一大群人聚集起来，几乎没有一个人留在家里；大家都惊叹他们的忍耐，因为他们被残忍地、非人地拖行。\par

总督折磨了宗徒斐理伯和他身边的圣人们之后，下令把他们带去，由神庙的祭司关在毒蛇偶像的神庙里，等他决定用哪种死刑来消灭他们每一个人。许多群众相信了基督的恩宠，归附了宗徒斐理伯和他身边的人，弃绝了毒蛇偶像，并因圣人们的坚忍而受到坚固，在信德中得以坚定；他们全体同声光荣天主，高呼阿们。\par

当他们——宗徒斐理伯、巴尔托羅买和玛利亚姆——被关在毒蛇神庙里时，毒蛇的祭司们聚集在同一地方，还有一大群群众，大约七千人；他们跑到总督那里，喊着说：\Quote{为我们向这些外国人、术士、败坏者和诱惑者报仇吧。因为他们自从来到我们这里，我们的城就充满了各种恶行；他们还杀死了我们女神的儿子——那些蛇；他们还封闭了神庙，祭台也荒废了；我们没有找到那原本拿来让毒蛇喝了好入睡的酒。但若你想知道他们真是术士，请看看他们如何想蛊惑我们，叫我们信了天主后活在贞洁和虔诚中；也看看他们如何进了城；又看看那些龙怎么没有弄瞎他们，甚至没有杀死他们；也看看他们怎么没有喝他们的血；但连那些保护我们城免受外人侵扰的，也被这些人推翻了。}\par

总督听了这些，越发怒火中烧，充满怒气与威胁；他极其愤怒，对祭司们说：\Quote{他们连我自己的妻子都蛊惑了，\Emph{你们还需要说什么？}从那时起，她就用奇怪的话对我说；她整夜祈祷，有光环绕，用奇特的语言说话；并且大声呻吟，说：\InnerQuote{耶稣真光已经临到我。}而我，从寝宫出来，想从窗子往下看，看看她说的那光的耶稣；它像闪电一样临到我，我几乎失明；从那时起，我就害怕我的妻子，因为她有那发光的耶稣。告诉我，你们这些祭司，我该怎么办？}他们对他说：\Quote{总督啊，我们确实不再算是祭司了；因为自从你把他们关进去，因着他们的祈祷，神庙不仅从根基摇动，而且肯定要倒塌了。}\par

于是总督下令，将斐理伯及同他在一起的人从庙中带出来，押解到审判台前，对刽子手说：\Quote{剥光斐理伯、巴尔多禄茂和玛利亚默的衣服，仔细搜查，看是否能找到他们的邪术。}因此，他们先剥了斐理伯的衣服，然后剥了巴尔多禄茂的，最后轮到玛利亚默；他们拖着她，说：\Quote{让我们把她剥光，让所有人都看见她如何跟随男人；因为她尤其欺骗了所有的妇女。}暴君对司祭们说：\Quote{在全城各处宣告，所有人，无论男女，都来看她的无耻行为，她和这些术士同行，无疑与他们通奸。}他下令将斐理伯吊起来，刺穿他的脚踝，并拿来铁钩，也刺穿他的脚跟，头朝下悬挂在庙对面的一棵树上；将巴尔多禄茂展开在斐理伯对面，把他的双手钉在庙门的墙上。\par

两人都笑了，彼此对视，斐理伯和巴尔多禄茂；因为他们仿佛并未受酷刑：因为他们的刑罚就是奖品和冠冕。当他们剥玛利亚默的衣服时，看，立刻她的身体形象在所有在场者面前发生了变化，立刻有一团火焰般的云围绕着她，出现在众人面前；他们再也无法看向圣玛利亚默所在的地方，全都逃离了她。\par

斐理伯用希伯来语对巴尔多禄茂说：\Quote{我们的兄弟若望在哪里？因为，看，我即将脱离身体；是谁为我们祈祷了？因为他们也违背情理，对我们的姊妹玛利亚默动了手；看，他们放火烧了斯塔基的家，说：让我们烧掉它，因为他接待了他们。那么，巴尔多禄茂，你愿意火从天上降下，把他们烧尽吗？}\par

斐理伯正这样说时，看，若望也进了城，像他们的一个同胞一样；他在街上走动，问：\Quote{这些人是谁？为什么受罚？}他们对他说：\Quote{你不可能本城的人，却问起这些人，他们伤害了许多人：因为他们关闭了我们的众神，用邪术砍断了蛇和龙；他们还复活了许多死人，使我们惊愕，猛烈地控诉我们，这些被悬挂的外乡人也想祈求天火，烧毁我们和我们的城。}\par

若望便说：\Quote{我们去吧，你们带我去看他们。}于是他们像对本城同胞一样，领若望到了斐理伯所在的地方；那里有一大群人，还有总督和司祭们。斐理伯看见若望，用希伯来语对巴尔多禄茂说：\Quote{兄弟，若望来了，他曾住在巴里克，那里有活水。}若望看见斐理伯头朝下被悬挂着，双脚脚踝和脚跟都被刺穿；他也看见巴尔多禄茂被展开在庙墙上；他对他们说：\Quote{那悬于天地之间者的奥迹必与你们同在。}\par

他又对那城的居民说：\Quote{你们这些住在\OriginalTerm{奥斐奥里玛耶辣颇里}{Ophioryma Hierapolis}的居民啊，你们当中何等巨大的无知！你们在错谬的路径上迷失了。龙的呼吸向你们吹气，使你们在三个方面盲目：即，使你们身体盲目，灵魂盲目，精神盲目；你们被毁灭者击中了。看看整个受造界，无论地上、天上还是水中，蛇与任何上界之物都无相似之处；它属于腐朽的种类，已被天主化归于无；因此它扭曲、弯曲，毫无生命；愤怒、暴怒、黑暗、火和烟充满它所有的肢体。那么现在，你们为什么惩罚这些人，仅仅因为他们告诉你们，蛇是你们的仇敌？}\par

他们听了若望这些话，就举手攻击他，说：\Quote{我们以为你是本城同胞，现在你却表明你是他们的同伴。像他们一样，你也要被处死；因为司祭们打算榨你的血，和酒混合，带去给毒蛇喝。}当司祭们试图抓住若望时，他们的手都麻痹了。若望对斐理伯说：\Quote{我们绝不可以恶报恶。}斐理伯对若望说：\Quote{看，现在，我的主耶稣在哪里？祂曾告诉我不要报复。但就我而言，我不能再忍受了；我要向他们实现我的威胁，将他们全部消灭。}\par

若望、巴尔托罗茂和玛利亚默阻止了他，说：我们的主被殴打、被鞭笞、被\Emph{钉在十字架上}、被灌以苦胆和酸醋，并说：\BibleQuote{父啊，宽赦他们吧，因为他们不知道他们做的是什么。}\BibleRef{路 23:34}祂也教导说：\BibleQuote{你们跟我学吧，因为我是良善心谦的。}\BibleRef{玛 11:29}因此，我们也应当忍耐。斐理伯说：\Quote{走开，不要软化我；因为我不能容忍他们把我头朝下吊起，用铁器刺穿我的脚踝和脚跟。而你，若望，天主所爱的人，你对他们讲了多少道理，他们都不听！所以，走开，我要诅咒他们，他们必被彻底毁灭，一个不留。}于是他开始诅咒他们，用希伯来语呼号：\OriginalTerm{阿巴罗，阿雷蒙，依杜撒耳，塔尔塞肋翁，纳科特，埃杜纳夫，特肋托洛伊}{Abalo, aremun, iduthael, tharseleon, nachoth, aidunaph, teletoloi}，即：\Quote{基督的父，唯一全能的天主！天主，万世敬畏你，大能而公正的审判者，在你权能中你的名是撒巴奥特，你当受赞美，直到永远。在你面前，天上的统治者和掌权者战栗，革鲁宾生灵那喷火的可怖势力也颤抖；你是圣洁威严的君王，你的名曾临于旷野的野兽，它们便被驯服，以理性的声音赞美你；你垂顾我们，乐意应允我们的祈求；你在我们被塑造之前就已认识我们；你是万物的监察者。现在，我恳求，让伟大的哈得斯张开它的口；让深渊吞没这些不敬虔的人，他们在这城中不愿意接受真理之言。愿如此成就，撒巴奥特。}看，深渊忽然裂开，总督所坐的整个地方、整个神庙、他们所崇拜的毒蛇、以及大群人和毒蛇的司祭——大约七千人——连妇女和儿童都被吞没，只有宗徒们所在之处安然无恙。总督被吞入深渊；他们的声音从地下哀哭而来：\Quote{天主的荣耀的宗徒们，求你怜悯我们，因为我们如今看到那些没有承认被钉者之人的审判：看，十字架光照了我们。耶稣基督，求你向我们显现，因为我们全都活着下到哈得斯，正在受鞭打，因为我们不义地钉死了你的宗徒。}有人听到一个声音说：\Quote{我将在光明之十字架上怜悯你们。}\par

留下来的人有斯塔希斯和他的全家、总督的妻子、以及与她一起信了主的五十个妇女，此外还有许多男女、一百个童女——因她们的贞洁未被吞没，且已受了基督的印记。\par

那时，主向斐理伯显现，说：\Quote{斐理伯，你没有听见吗：\BibleQuote{不可以恶报恶}？你为什么降下这样的毁灭？斐理伯，谁把手扶在犁上，却往后看，\BibleRef{路 9:62}他的犁沟能正直吗？谁把自己的灯交给别人，自己却坐在黑暗里？谁舍弃自己的住所，自己却住在粪堆上？谁在冬天施舍自己的外衣，自己却赤身露体？有什么仇敌因恨他的人快乐而欢喜？有什么士兵不带全副武装就去打仗？有什么仆人完成主人的命令后不得称赞？谁在赛场上奋力奔跑却不接受奖赏？谁洗净了自己的衣服却故意弄脏它们？\BibleQuote{看，我的婚房已备好；但唯有那穿着发光衣服被找到的人有福}\BibleRef{玛 22:11}他才是头上接受冠冕的人。看，晚宴已备好；被邀请且准备好去赴宴的人有福了。\BibleQuote{田里的庄稼很多}\BibleRef{玛 9:37}好的工人有福了。请看百合花和一切花朵，好的农夫是首先享有它们的人。斐理伯，你怎么变得如此不仁慈，在愤怒中诅咒你的仇敌？}\par

斐理伯说：\Quote{主，你为什么向我发怒，因我诅咒了我的仇敌？你为什么还不把他们踩在脚下，因为他们还活在深渊里？主，你知道我因你的缘故来到这城中，并因你的名追查了所有偶像的错误和所有魔鬼。龙和蛇都已枯萎。既然这些人不接受你的光，我就诅咒了他们，他们就活着下到哈得斯去了。}\par

救主对斐理伯说：\Quote{但因你违背了我，以恶报恶，没有遵守我的诫命，为此你将要光荣地完成你的赛程，并由我的圣天使引领，与他们一同到达喜乐的乐园；他们确实要与我一同进入乐园，但我将命令你被关在乐园之外四十天，在旋转的火焰之剑的恐怖之下，你将叹息，因为你加害了那些加害你的人。四十天后，我将派遣我的总领天使弥额尔，他握住守护乐园的剑，将你领进乐园，你将看见所有那些以无玷之身行走的义人，而后你将在天上朝拜我父的荣耀。然而，你离世的记号将在我的十字架上得荣耀。巴尔托罗茂将前往吕考尼雅，在那里他也要被钉十字架；玛利亚默将把她的身体安放在约但河中。但我，斐理伯，将不容忍你，因为你把这些人吞入了深渊；但是，看，我的圣神在他们里面，我将从死人中把他们提拔上来；这样，当他们看见你时，就将相信那派遣你者的荣耀。}\par

救主转过身来，举起手，在空中画了一个十字，从上直到深渊，那十字充满了光，形状如同梯子。所有从城里下到深渊的人群都上了这光明的十字之梯；但总督和他们所敬拜的毒蛇仍留在下面。当人群上来后，看见斐理伯头朝下悬挂着，便为所行的不法行为大声哀哭。他们也看到了巴尔多禄茂和玛利亚姆，她恢复了原来的容貌。看，主在斐理伯、巴尔多禄茂、玛利亚姆、斯塔基以及所有不信的民众眼前升上了天，他们默不作声，恐惧战兢地光荣天主。所有民众都喊说：\Quote{唯独他是天主，这些人真实地传扬他；唯独他是天主，派遣这些人来为我们的救恩。让我们真正为我们的大错悔改，因为我们绝不配得永生。现在我们相信了，因为我们看见了伟大的奇迹，因为救主将我们从深渊中带了上来。}他们都俯伏在地，朝拜斐理伯，并恳求他（他正要离开）：\Quote{不要再行一个奇迹，又把我们送回深渊。}他们祈祷，使自己堪当见到基督的显现。\par

斐理伯仍然悬挂着，对他们说：\Quote{你们要听，要明白我的天主的能力何其伟大，要记住你们在下面所看见的，以及你们的城如何被倾覆，只有接待我的那房屋幸免；现在我的天主的甘甜已将你们从深渊中领上来，而为了你们的缘故，我必须绕行乐园四十天，因为我曾向你们发怒，要报复你们。只有这条诫命我没有遵守，就是我没有以善报恶。但我告诉你们：从今以后，因天主的良善，弃绝邪恶，好使你们堪当承受主的感恩祭。}\par

有些信众跑上去要把斐理伯取下来，除去他脚踝上的铁钩和钩子。但斐理伯说：\Quote{我的孩子们，不要这样做，不要为此接近我，因为我的结局就是这样。你们这些在主内蒙光照的人，请听我说：我来到这城，不是为了做买卖或做别的事；而是我注定要在这城中，以你们所见的这种情况离开我的身体。所以，不要因我这样悬挂而悲伤；因为我带有第一个人的印记，他头朝下被带到地上，又通过十字架的木头，从违命的死亡中复生。现在我完成所吩咐我的事；因为主对我说：\InnerQuote{除非你使那属下的成为属上的，使那属左的成为属右的，你绝不能进入我的国。}所以不要像那不变的模型，因为整个世界都改变了，每个住在身体内的灵魂都忘记天上的事；但我们这些拥有天上光荣的人，不要寻求那外在的，即身体和奴隶的房屋。不要不信，但要相信，并彼此宽恕过失。看，我悬挂了六天，我从真正的审判者那里受到责备，因为我竟然以恶报恶，在我的正直之路上设置了绊脚石。现在我要升到高处；不要忧愁，反而要喜乐，因为我正在离开这住所，我的身体，逃脱了那惩罚每个在罪中的灵魂的毒龙。}\par

斐理伯环视群众，说：\Quote{你们从死者中从阴府和深渊的吞没中上来的人——那光明的十字架因圣父、圣子、圣神的良善将你们领上高处——那位天主成为人，由童贞玛利亚取得肉身，是不朽的，住在肉身内；他死了，又复活了死者，怜悯了人类，除去了罪的毒钩。他是伟大的，却为了我们成为渺小的，直到他使渺小的扩大，并领他们进入他的伟大。他拥有甘甜；他们却唾弃他，给他喝苦胆，为要使那些对他怀有苦毒的人尝到他的甘甜。所以你们要依附他，不要离弃他，因为他是我们永远的生命。}\par

\Person{斐理伯}一说完这项宣告，便向他们说：\Quote{解开\Person{巴尔多禄茂}；}他们上前解开了他。解开后，\Person{斐理伯}向他说：\Quote{\Person{巴尔多禄茂}，我在主内的弟兄，你知道主派你同我来到这城，你与我们姊妹\Person{玛利亚默}一起分担了一切危险；但我知道你离世被定在\Place{吕考尼亚}，\Person{玛利亚默}被定在\Place{约旦河}离世。所以现在我命令你：等我离世后，你要在这地方建一座教堂；让豹和山羊羔进入教堂，作为信者的记号；让\Person{尼卡诺拉}供养他们，直到他们离世；他们离世后，把他们葬在教堂门口。将你的平安留在\Person{斯塔希}的家，如同基督将祂的平安留在这城一样。让所有信主的童贞女每天在那家中，两人一组照顾病人；但不可与青年男子来往，免得撒殚诱惑你们：\BibleRef{格前 7:5}因为他是爬行的蛇，曾藉厄娃使亚当滑入死亡。愿此时不再发生如同厄娃的事。但是，\Person{巴尔多禄茂}，你要好好照顾她们；你要将这些吩咐交给\Person{斯塔希}，并立他为主教。不可把主教职位交给年轻人，免得基督的福音受辱；凡教导人的，应有与言语相称的行为。我现在要到主那里去，取了我的身体，用叙利亚纸张预备埋葬；不要在我周围放亚麻布，因为我主的身体是用细麻布包裹的。用纸莎草茎将我的身体捆绑结实，葬在教堂里；为我祈祷四十天，使主赦免我报复那恶待我之人的过犯。看，\Person{巴尔多禄茂}，我的血滴在地上的地方，将从我的血中长出一棵植物，成为葡萄树，结出一串葡萄；取那串葡萄，挤入杯中；第三天领受后，高呼阿们，使祭献得以完成。}\par

\Person{斐理伯}说完这话，便这样祈祷：\Quote{主耶稣基督，万世之父，光明之王，祢以祢的智慧使我们有智慧，赐给我们祢的悟性，并把祢良善的计谋赐予我们，祢从未在任何时候离弃我们；祢是那除去投奔祢者疾病的那一位；祢是永生天主之子，祢赐给我们祢智慧的临在，祢行神迹奇事，使迷途者回转；祢给战胜仇敌者加冕，祢是卓越的审判者。现在来吧，耶稣，赐给我战胜一切敌对权势的永恒胜利之冠，当我渡过火的水和所有深渊时，不要让他们的黑暗空气遮蔽我。我主耶稣基督，不要让仇敌在祢的审判座前有控告我的理由；但给我穿上祢荣耀的长袍，祢那永远闪耀的光明印记，直到我经过世上一切权势和那埋伏等候我们的恶龙。所以，我主耶稣基督，使我在空中与祢相遇，赦免了我对我仇敌所施的报复；将我的身体形状转变为天使的光荣，使我在祢的福乐中安息；让我从祢那里领受祢曾应许给祢圣人们的永远承诺。}\par

\Person{斐理伯}说完这话，便断了气，众百姓都看着他，哭泣着说：\Quote{这灵魂的生命在平安中完成了。}他们便说：阿们。\par

\Person{巴尔多禄茂}和\Person{玛利亚默}取下他的身体，照\Person{斐理伯}所吩咐的做了，葬在那地方。立刻有声音从天上来：\Quote{宗徒\Person{斐理伯}已被耶稣基督、竞赛的审判者，用不朽的冠冕加冕。}众人都高呼：阿们。\par

三天后，圣\Person{斐理伯}的血滴落的地方长出了葡萄树的植物。他们遵从他的一切命令，连续四十天献祭，不住祈祷。他们在那里建了教堂，立\Person{斯塔希}为教堂的主教。\Person{尼卡诺拉}和所有信友聚集，不住地光荣天主，因他们中间所发生的奇迹。全城的人都相信了耶稣的名。\Person{巴尔多禄茂}命令\Person{斯塔希}为那些相信的人授洗，因父、及子、及圣神之名。\par

四十天后，救主以\Person{斐理伯}的形像显现，对\Person{巴尔多禄茂}和\Person{玛利亚默}说：\Quote{我所爱的弟兄们，你们愿意在天主的安息中安息吗？乐园已为我敞开，我进入了耶稣的光荣。去吧，到为你们指定的地方去；因为在这城中被分别并栽种的植物必结出丰盛的果实。}于是他们向弟兄们致候，并为他们每人祈祷后，离开了\Place{奥斐奥里玛}城，即\Place{亚细亚}的\Place{希拉波利斯}；\Person{巴尔多禄茂}往\Place{吕考尼亚}去了，\Person{玛利亚默}往\Place{约旦河}去了；\Person{斯塔希}和同他在一起的人留下，维持着教会，在基督耶稣我们的主内，愿光荣与权能归于祂，直到永远。阿们。\par

% structure: p0039 source=h2 target=section reason=relative-heading-rank; base=h2

\section{圣宗徒\Person{斐理伯}前往\Place{上希腊}时的事迹}

那时，\Person{斐理伯}进入名为\Place{希腊}的\Place{雅典}城，有三百位哲学家聚集到他那里，说：\Quote{我们去看看他的智慧如何；}因为关于\Place{亚细亚}的智者，他们说他们的智慧很伟大。他们以为\Person{斐理伯}是哲学家，因为他穿着隐修者的衣服；他们不知道他是基督的宗徒。因为耶稣赐给祂门徒的服装只是一件外衣和一块亚麻布。\Person{斐理伯}就是这样行走的。因此，\Place{希腊}的哲学家看见他时，便害怕了。于是他们聚集在一处，彼此说：\Quote{来，让我们查阅我们的书籍，免得这个外乡人胜过我们，使我们蒙羞。}\par

他们这样做了以后，便聚集到同一处，对斐理伯说：\Quote{我们有祖传的教义，我们以此为乐，追求知识；但若你有什么新事，陌生人，请毫不嫉妒地大胆展示给我们；因为我们别无所需，只想听一些新事。}\BibleRef{宗 17:21}\par

斐理伯回答他们说：\Quote{希腊的哲学家们，如果你们愿意听一些新事，并渴望新事，就应当抛弃旧人的性情，如同我主所说：\InnerQuote{不能把新酒装在旧皮囊里，以免皮囊破裂，酒漏出来，皮囊也坏了。而是把新酒装在新皮囊里，这样两者都能保全。}这些事是主用比喻说的，以祂神圣的智慧教导我们，许多人会喜爱新酒，却没有新鲜的新皮囊。我爱你们，希腊人，我祝贺你们说了：\InnerQuote{我们喜爱新事。}因为我的主确实把崭新而清新的教导带到了世界上，为要扫除世上的一切教导。}\par

哲学家们说：\Quote{你称谁为主？}\newline{}斐理伯说：\Quote{我的主就是在天上的耶稣。}\newline{}他们又对他说：\Quote{请毫不嫉妒地将祂显明给我们的悟性，好使我们也能相信祂。}\newline{}斐理伯说：\Quote{我将要介绍给你们的那位主，超乎万名之上；再没有别的名。}\BibleRef{弗 1:21}\Quote{我仅说这一点：正如你们所说的——不要因嫉妒而拒绝我们——愿我不拒绝你们；反而我要极其欢欣踊跃地向你们启示那个名，因为我在世界上没有别的工作，只有这一宣告。因为我的主来到这世界时，祂拣选了我们十二人，以圣神充满了我们；从祂的光中，祂使我们知道祂是谁，并命令我们藉着祂宣讲一切救恩，因为天下没有赐下别的名，使我们赖以得救。}\BibleRef{宗 4:12}\Quote{为此，我来到你们这里，要使你们完全确信，不仅藉言语，也藉着在我主耶稣基督的名内所行的奇妙作为。}\par

哲学家们听了这话，就对斐理伯说：\Quote{我们从你那里听到的这个名，从未在我们祖传的书中见过；因此，现在我们怎能知道你的话呢？}此外，他们又对他说：\Quote{请宽限我们三天，让我们彼此商议这个名；因为我们对这事并非不重视——背弃我们祖传的宗教。}斐理伯于是对他们说：\Quote{你们随意商议吧；因为此事毫无诡诈。}\par

那三百位哲学家集合起来，彼此交谈说：\Quote{你们知道这个人带来了一种奇怪的哲学，他所说的话使我们迷惑。那么，我们该怎样对待他，或者对待他所讲的那位被称为耶稣、万世君王的名呢？}此外，他们又彼此说：\Quote{我们确实无法与他辩论，但犹太人的大司祭可以。因此，若以为好，就派人去请他，好让他来与这个陌生人辩论，以便我们准确了解所宣讲的名。}\par

他们就这样写信给耶路撒冷：——\newline{}\Quote{希腊的哲学家们致在耶路撒冷的犹太人大司祭\Person{阿纳尼亚斯}。我们之间始终存在深厚……如你所知，我们雅典人是寻求真理者。有一个外乡人来到希腊，名叫\Person{斐理伯}；简言之，他以言语和非凡的奇迹极大地扰乱了我们，他引入了一个荣耀的名——耶稣，并自称是祂的门徒。他还行了许多奇事，我们写信告诉你：他驱逐了久附人身的多魔，使聋子听见，瞎子看见；更奇妙的是——这也是我们应该首先提及的——他使已经死去了相当天数的人复活。他的名声传遍了整个希腊和马其顿；从周围城市有许多人来到他那里，带着患有各种疾病的人，他都藉着耶稣的名治愈了他们。因此，请你毫不迟疑地到我们这里来，好由你亲自向我们宣告，耶稣——他所教导的这个名——到底意味着什么。正是为此，我们写了这封信给你，大司祭。}\par

他读完后，充满了极大的愤怒，撕裂了衣服，说：\Quote{那个骗子竟然也到雅典，在哲学家中间，迷惑他们吗？}\OriginalTerm{曼塞玛特}{Mansemat}——就是撒殚——悄然进入了阿纳尼亚斯，使他充满了愤怒和狂怒；他说：\Quote{如果我让斐理伯本人和与他在一起的人活着，法律就会被完全摧毁，他们的教导可能会充满全地。}大司祭回到自己家中，经师和法利塞人也来了；他们彼此商议说：\Quote{我们该怎样应对这些事？}他们对大司祭阿纳尼亚斯说：\Quote{起来，武装自己，带领五百名精壮百姓，前往雅典，务必杀死斐理伯，这样你就能推翻他的教导。}\par

他穿上大司祭的礼服，带着那五百人，以极大的排场来到希腊。斐理伯正在城中一位首领的家里，与已经信主的弟兄们在一起。大司祭和随从，以及那三百位哲学家，来到斐理伯所在的房屋门口；有人告诉斐理伯说他们在外面。他就起来出去了。大司祭看见他，就对他说：\Quote{哦，斐理伯，术士和魔法师，因为我认识你，在耶路撒冷，你的师傅那个骗子曾称你为雷霆之子。难道整个犹太地还不够你活动，你竟也来到这里，欺骗寻求智慧的人吗？}斐理伯说：\Quote{但愿，阿纳尼亚斯，你心中的不信之幕被除去，这样你就能明白我的话，并从中知道究竟是我还是你才是骗子！}\par

\Person{阿纳尼雅}听了这话，对\Person{斐理伯}说：\Quote{我要回答一切。}\Person{斐理伯}说：\Quote{你说吧。}大司祭说：\Quote{希腊人啊，这个\Person{斐理伯}相信一个名叫\Person{耶稣}的人，他生在我们中间，教导这异端，毁坏了法律和圣殿，使通过\Person{梅瑟}的净化和新月归于无有，因为他说：\InnerQuote{这些不是天主所命令的。}当我们看到他这样毁坏法律时，就起来反对他，将他钉在十字架上，使他的教导不得成就。因为他带来了许多改变；他作了恶见证，因为他像外邦人一样，共通享用一切并与血掺杂。我们将他交出，处死他，将他埋葬在坟墓里；他的门徒偷走了他，到处宣扬他从死者中复活了，并宣称他在天上坐在天主的右边，使许多人走入歧途。但现在这些人，自己和我们一样有割损，却不遵守它，因为他们开始在\Place{耶路撒冷}藉\Person{耶稣}的名行许多大能的作为；他们被逐出\Place{耶路撒冷}后，周游世界，用那个\Person{耶稣}的魔术欺骗所有人，正如现在这个\Person{斐理伯}也来用同样的手段欺骗你们。但我将把他带回\Place{耶路撒冷}，因为\Person{阿赫劳斯}王也在搜寻他，要杀他。}\par

当周围的群众听到这话时，那些在信德上坚固的人并没有动摇或犹疑，因为他们知道\Person{斐理伯}要在\Person{耶稣}的光荣中得胜。因此\Person{斐理伯}以基督的德能，非常大胆地、欢跃地陈词说：\Quote{我，\Place{雅典}人啊，以及你们中的哲学家们，来到你们这里，不是要用言语教导你们，而是藉着显示奇迹；部分地你们已经迅速看见了藉着我所成就的事，藉着那个名字，大司祭自己已被弃绝。因为，看，我要呼求我的天主，教导你们，而你们将证实双方的话。}\par

大司祭听了这话，跑到\Person{斐理伯}那里，想要鞭打他，就在那一刻，他的整只手枯干了，眼睛也瞎了；同样，与他在一起的五百人也瞎了眼。那五百人就辱骂并诅咒大司祭，说：\Quote{我们从\Place{耶路撒冷}出来时对你说过：\InnerQuote{住手吧}；因为我们是人，不能与天主争斗。但我们恳求你，\Person{斐理伯}，天主\Person{耶稣}的宗徒，赐给我们藉着他而来的光明，使我们也能真正成为他的奴隶。}\par

\Person{斐理伯}看到所发生的事，就说：\Quote{软弱的本性啊！你冲向了我，却立刻降卑；苦涩的海啊！你兴起波澜反对我，想要将我们抛出去，却自己使波浪平息。所以如今，我们良善的管家\Person{耶稣}，圣光啊，你没有忽略我们这些一起以各种善工向你呼求的人，反而来藉着我们完成它们。所以现在，主\Person{耶稣}，来吧，指正这些人的愚妄。}\par

大司祭对\Person{斐理伯}说：\Quote{你难道想要使我们背离祖先的传统、旷野的天主和\Person{梅瑟}吗？你以为你会使我们成为纳匝肋人\Person{耶稣}的追随者吗？}然后\Person{斐理伯}对他说：\Quote{看，我要呼求我的天主来临，在他面前显现给你和这五百人，以及这里所有的人；或许你会改变心意而相信。但即使你至终仍停留在不信中，一件奇特的事要临到你，这事要被世代相传——你也要活着下到阴府，在所有看见你的人面前，因为你仍留在不信中，也因为你想使这群人远离真生命。}于是\Person{斐理伯}祈祷说：\Quote{圣子的圣父，\Person{耶稣基督}啊，你恩准了我相信他，请派遣你的爱子\Person{耶稣基督}来临，指正这不相信的大司祭，使你的名在爱子\Person{基督}内受光荣。}\par

当\Person{斐理伯}还在呼喊这些话时，忽然天开了，\Person{耶稣}以极其光荣的荣耀，带着闪电降临；他的面容比太阳光耀七倍，他的衣服比雪更白，因此\Place{雅典}所有的偶像都突然倒在地上。众人在痛苦中逃跑；住在他们当中的魔鬼喊说：\Quote{看，我们也因为那已向城市显现的、天主子\Person{耶稣}而逃跑。}于是\Person{斐理伯}对大司祭说：\Quote{你听见魔鬼因那被看见的那位而喊叫，你却不相信那临在的那位是万有的主吗？}大司祭说：\Quote{除旷野中的那位以外，我没有别的天主。}\par

正当\Person{耶稣}升天时，发生了极强烈的地震，以致他们站的地方裂开了；群众跑过来，扑倒在宗徒脚前，喊说：\Quote{天主的仆人，可怜我们吧！}同样，那五百人也再次喊道：\Quote{\Person{斐理伯}，可怜我们吧，使我们能认识你，并藉着你认识生命之光\Person{耶稣}；因为我们曾对这不相信的大司祭说：\InnerQuote{我们是有罪的人，不能与天主争斗。}}\par

于是斐理伯说：\Quote{我们里面没有仇恨，但基督的恩宠将使你们恢复视力；然而我要先使大司祭恢复视力，好使你们因此更加相信。}有一个声音从天上对斐理伯说：\Quote{斐理伯，昔日是雷霆之子，如今却是温良之子，无论你向我父求什么，他必为你成就。}众百姓听见那声音，都惊慌失措，因为那声音比雷声还大。斐理伯便对大司祭说：\Quote{因我主之声的能力，阿纳尼雅，恢复视力吧！}他立刻就看见了，环顾四周，说：\Quote{耶稣的邪术有什么了不起，这斐理伯片刻之间使我失明，又片刻之间使我复明？}斐理伯说：\Quote{那么你相信耶稣吗？}大司祭说：\Quote{你以为你能迷惑我，说服我吗？}与他同来的五百人，听见他们的大司祭虽然复明却仍不信，就对旁边的人说，求斐理伯使他们也恢复视力，\Emph{他们说}：\Quote{好叫我们除掉这个不信的司祭长。}\par

斐理伯说：\Quote{不要对恶人报复。}他对大司祭说：\Quote{你身上将有某个重大的记号。}大司祭对斐理伯说：\Quote{我知道你是个术士，是耶稣的门徒；你迷惑不了我。}宗徒便对耶稣说：\OriginalTerm{萨巴坦、萨巴塔布、布拉马努赫}{Sabarthan, sabathabt, bramanuch}，快来。随即，阿纳尼雅所在之处的地裂开，吞没他至膝。阿纳尼雅喊道：\Quote{哎，这真巫术的力量真大，因斐理伯用希伯来语威胁并命令地，地就裂开；它抓住我直到膝盖，脚后跟好像有些钩子往下拉我，好叫我信斐理伯；但他不能说服我，因为在耶路撒冷我就知道他的幻术。}\par

斐理伯发怒说：\Quote{地，你紧紧抓住他，直到肚脐。}地立即将他拉下去。他说：\Quote{我下面的一只脚变成冰，另一只却烫得可怕；但你的邪术，斐理伯，我决不屈服。除非我下面剧痛难忍，我根本不信。}群众想要用石头砸他。斐理伯说：\Quote{不可这样；这事暂时发生，他被吞没至肚脐，是为使你们的灵魂得救，因为他几乎用恶言把你们拖入不信。但即使他悔改，我也会把他从地里拉上来，为拯救他的灵魂；然而他确实不配得救。那么，若他继续不信，你们要看见他沉入深渊，除非主愿意使阴府中的人复活，叫他们承认耶稣是主。因为在那一日，万口必承认耶稣是主\BibleRef{斐 2:11}，并且圣父与圣子有同一光荣，偕同圣神，至于永远。}\par

斐理伯说完这话，伸出右手，在空中划向那五百人，因耶稣之名。他们的眼睛就开了，全都异口同声赞美天主说：\Quote{我们赞颂你，基督耶稣，斐理伯的天主，因你驱除了我们的盲目，赐给我们你的光明，即福音。}斐理伯因他们的话极其欢喜，因为他们如此在信德中坚定了。此后，斐理伯转身对大司祭说：\Quote{你也要以纯洁的心承认耶稣是主，好使你得救，像你身边的人一样。}但大司祭嘲笑斐理伯，仍然不信。\par

斐理伯见他仍不信，就看着他，对地说：\Quote{张开你的口，在那些信了基督耶稣的人面前，将他吞没直至颈项。}在同一时刻，地张开了口，吞没他直至颈项。群众因所发生的奇事彼此议论。\par

城中有一位首领跑来，喊着说：\Quote{有福的宗徒，有个恶魔攻击了我的儿子，对我喊说：\InnerQuote{你既让一个外人进城，你又是首先废除我们敬拜和祭献的人，我除了杀死你这独生子，还能为你做什么？}他说完这话，就勒死了我的儿子。所以现今我恳求你，基督的宗徒，不要使我的喜乐变为哀伤，因为我也相信了你的话。}\par

宗徒听了这话，说：\Quote{我惊奇魔鬼的活动，它无所不在，竟敢攻击那些我未能前去帮助的人，正如他们现在试探你们，想要使你们跌倒。}他就对那人说：\Quote{将你的儿子带来，我必藉着我的基督，将他活着交还给你。}那人欢喜地跑去带他的儿子。当他走近自己的房子时，喊道：\Quote{我的儿子，我来接你到宗徒那里，使他将你活着交还给我。}他吩咐仆人抬床；他的儿子二十三岁。斐理伯看见他，动了心；他转向大司祭说：\Quote{这事的发生是为给你一个机会；那么，若我使他复活，你从此就相信吗？}大司祭说：\Quote{我知道你的邪术，你会使他复活，但我决不信你。}斐理伯发怒说：\Quote{你该受诅咒！那么，在众人面前，完全下到深渊去吧！}在同一时刻，他活着下了阴府，只是大司祭的长袍从他身上飞脱；因此，从那天起，没有人知道那司祭长袍的下落。宗徒转身为那男孩祈祷；赶走魔鬼后，使他复活，将他活着置于他父亲身边。\par

众人看见这事，就喊说：\Quote{\Person{斐理伯}的天主是唯一的天主，祂惩罚了大司祭的不信，赶走了少年身上的魔鬼，并使他从死者中复活。}那五百人看见大司祭被吞入深渊，以及其他的奇迹，就恳求\Person{斐理伯}，\Person{斐理伯}就将基督的印记赐给他们。\Person{斐理伯}在\Place{雅典}住了两年；建立了一座教会后，任命了一位主教和一位司铎，然后就离开前往\Place{帕提亚}，宣讲基督。愿光荣归于祂，直到永远。阿们。\par

% structure: p0064 source=h2 target=section reason=relative-heading-rank; base=h2

\section{\WorkTitle{斐理伯行传}补篇}

（来自巴黎手抄本）\par

他这样教导他们说：\Quote{我的弟兄们，我圣父的儿女——因为你们在基督内属于我的家族，是我城邑（天上的\Place{耶路撒冷}）的本体，我居所的喜乐——你们为什么被你们的仇敌，那扭曲、弯曲、邪僻的蛇所掳掠？天主没有给它手和脚。它的行走是弯曲的，因为它是恶者的儿子；它的父亲是死亡，母亲是败坏，毁灭在它的身体内。那么，不要进入它的毁灭；因为你们被它儿子的不信和欺骗所束缚，它儿子是无秩序的，没有本体，无形状，在整个受造界中没有形状，无论是在天上、地上，或水中的鱼中。但如果你们看见它，要躲避它，因为它不像人；它的居所是深渊，它行走在黑暗中。因此要躲避它，免得它的毒液倾泻在你们身上；如果它的毒液倾泻在你们身上，你们就行走在它的邪恶中。你们要停留在真崇拜中，做忠信、虔敬、良善、无诡诈的人。要躲避\Person{撒殚}这龙，从你们身上除掉它的邪恶种子，即欲望，欲望在灵魂中生出疾病，就是蛇的毒液。因为欲望从起初就是属于蛇的，它武装自己攻击忠信的人；因为它是从黑暗中出来，又回到黑暗中。因此，你们来到我们这里，或者更确切地说，通过我们来到天主面前，就应该从你们身上扔掉魔鬼的毒液。}\par

当宗徒正说这些话时，看，\Person{尼卡诺拉}从她的房子里出来，和她的奴隶们一起进入\Person{斯塔基斯}的家。当她走近房门口时，看，\Person{玛利安}用叙利亚语对她说：\Quote{Helikomaei, kosma, etaa, mariacha}。然后她解释她的话说：\Quote{啊，圣神的女儿，你是我的主母，你曾被许给蛇作为抵押；但我来是要解救你：我要折断你的捆绑，从根部切断它们。看，解救你的释放者已经来到；看，正义的太阳已经升起，要光照你。}\par

当她这样说时，那阴郁的暴君跑来，气喘吁吁。在门前的\Person{尼卡诺拉}听见了，就在众人面前鼓起勇气，喊说：\Quote{我是希伯来人，是希伯来人的女儿；用我祖先的语言对我说话，因为我已听见你们的宣讲，并已从这疾病中痊愈。我尊敬并光荣天主的良善，因祂使你们在这地上完全被剥夺。}\par

当她这样说时，暴君来了，抓住她的衣服，说：\Quote{啊，\Person{尼卡诺拉}，我不是让你因疾病躺在床上吗？你从哪里找到这能力和力气，能来到这些魔术师这里？除非你告诉我谁是医治者，我要严厉惩罚你。}\Person{尼卡诺拉}回答说：\Quote{啊，暴君的养育者，从你自己身上抛弃这暴虐，忘记你的邪恶行为，放弃这暂时的生命，除去虚浮的荣耀，因为它像影子一样过去；要寻求那永久的，从你自己身上除掉兽性而不敬虔的卑劣欲望的工作，拒绝虚浮的交往，那是死亡的耕种，黑暗的监狱；推倒败坏的中间隔墙，为自己预备贞洁无瑕的生活，好让我们都能生活在圣洁中。如果你愿意我与你一同留下，我将与你一同节制地生活。}\par

暴君听了这些话，就抓住她的头发，拖着她，用脚踢她，说：\Quote{你被我的刀杀死，比被这些外国魔术师和欺骗者看见更好。因此，我要惩罚你，并将欺骗你的人处死。}他愤怒地转向跟随他的刽子手，说：\Quote{把这些骗子带到我这里来。}刽子手跑到\Person{斯塔基斯}的家，抓住\Person{斐理伯}、\Person{巴尔托罗茂}和\Person{玛利安}，连同豹和山羊羔，拖着他们带来了。\par

暴君看见了他们，就向他们咬牙切齿，说：\Quote{把这些魔术师和欺骗者拖走，他们说要崇拜天主，欺骗了许多妇女的灵魂。}他命人带来皮带，绑住他们的脚。他下令从城门一直拖到神庙。许多群众聚集到那地方。他们对豹和山羊羔极其惊奇，因为它们像人一样说话；有些群众相信了宗徒们的话。\par

司祭们对暴君说：\Quote{这些人是魔术师。}他听了，就怒火中烧，充满气愤；他下令剥去\Person{斐理伯}、\Person{巴尔托罗茂}和\Person{玛利安}的衣服，说：\Quote{搜查他们。也许你们会找到他们的巫术。}刽子手剥了他们的衣服，抓住\Person{玛利安}，拖着她，说：\Quote{揭露她，好让他们知道是跟随他们的女人。}他下令拿棍子和坚固的绳索；穿过\Person{斐理伯}的脚踝后，他们拿来钩子，把绳索穿过他的脚踝，把他头朝下挂在一棵树前，那树在神庙门前；他们又把木桩钉进神庙墙里，就离开了。他们又捆绑\Person{巴尔托罗茂}的手脚，将他赤裸地伸展在墙上；剥去\Person{玛利安}的衣服时，她的身体外表变了，成为一个充满光的玻璃匣，他们无法靠近她。\par

斐理伯用希伯来语对巴尔多禄茂说：\Quote{在我们需要的今日，若望在哪里？因为看，我们正从身体中被交出。他们过分地动手对付玛利亚姆，又鞭打了豹子和山羊羔，并放火烧了斯塔希的房子，因为他接待了我们。因此，让我们说话，使火从天降下烧灭他们。}\par

斐理伯正这样说时，看，若望进了城，在街上行走，问城中的人说：\Quote{这骚动是什么？这些人是谁？为什么他们受刑罚？}他们对他说：\Quote{你不是这城的人吗？你不知道这些人吗？他们如何搅乱我们的房屋和全城？此外，他们甚至以宗教为借口说服我们的妻子离开我们，宣扬一个陌生的名字，即基督的名字；他们还用他们所拥有的巫术关闭了我们的神庙，并用我们从未知晓的外邦名字杀死了城中的蛇。他们住在盲人斯塔希的家中，那女人用唾沫使他恢复了视力；也许她就是所有巫术的源头。有一头豹子和一只山羊羔跟着他们，像人一样说话。如果你曾见过这样的事，你就不会被他们困扰。}若望回答他们说：\Quote{带我去见他们。}他们便带他到斐理伯被悬挂的庙宇。斐理伯看见若望，就对巴尔多禄茂说：\Quote{我的兄弟，看，\OriginalTerm{巴勒加}{Barega}的儿子——即活水——来了。}若望看见斐理伯头朝下被脚踝绑着悬挂，也看见巴尔多禄茂被绑在庙墙上。\par

他对城中的人说：\Quote{蛇的子孙哪，你们的愚昧何等大！因为欺诈之道欺骗了你们，恶龙的气息吹在了你们身上：你们为什么因这些人说蛇是你们的仇敌而惩罚他们？}\par

他们听了若望这些话，就下手抓他，说：\Quote{我们曾称你为同城公民，但现在你的言语显明你也与他们相通。因此，你也要受和他们一样的死，因为司祭们已这样决定：我们让他们头朝下悬挂时放干他们的血，与酒调和，献给毒蛇。}\par

他们正这样说话时，看，玛利亚姆从她所在的地方站起来，恢复了她先前的容貌。\par

司祭们向若望伸手，想要抓住他，却不能。于是斐理伯与巴尔多禄茂对若望说：\Quote{耶稣在哪里？祂命令我们不可亲手报复折磨我们的人。因为此后我必不再容忍他们。}斐理伯用希伯来语说：\Quote{我的父\OriginalTerm{乌撒尔}{Uthael}，即基督，威严之父，万世畏惧其名，全能者，宇宙的能力，其名以主权发出，\OriginalTerm{厄罗阿}{Eloa}：你永受赞美；诸统权者和权势者畏惧你，在你面前战栗；荣耀的君王！威严之父！你的名已传到旷野的野兽，它们因你而安静，并通过你，蛇已离开我们：在我们祈求前就听我们。你是在我们呼求前就看见我们的，你知道我们的思念，万物的监察者，从自身发出无数慈悲者；愿深渊开口，吞没这些不承认你真理之言的无神之人。}\par

就在那时，深渊开口，那地方猛烈震动，从总督到众百姓连同司祭们，他们都沉下去了。宗徒们和所有与他们在一起的人所在之处，以及斯塔希的房子、暴君的妻子尼加诺拉、从丈夫那里逃出来的二十四位妻子、以及未曾认识男子的四十位贞女，这些地方都未受震动。只有这些人没有下到深渊，因为他们已成为仆人，并接受了天主的话语和祂的印号；但城中其余所有人都被吞下到深渊里了。\par

救主在那时刻显现，对斐理伯说：\Quote{谁把手扶在犁上，却转身不使犁沟笔直？或是谁将自己的光给别人，自己却坐在黑暗中？或是谁住在污秽中，却将自己的居所留给外人？或是谁脱下衣服，在冬日赤身外出？或是哪个奴仆做了主人的事，却不被主人叫来赴宴？或是谁在赛场上热心奔跑，却得不到奖赏？斐理伯，看，我的婚房已预备好了，那有自己发光衣服的人是有福的；因为这人将喜乐的冠冕戴在头上。看，宴席已预备好了，被新郎邀请的人是有福的。田地的庄稼很大，能干的工人是有福的。}\par

斐理伯听了救主这些话，回答他说：\Quote{你曾允许我们，纳匝肋的耶稣，难道你不命令我们击打那些不愿你统治他们的人吗？但我们知道，你的名尚未传遍全世界，你差遣我们到这城来。我本不打算进入这城，你赐给我你真实的诫命后，差遣了我，要我赶走一切欺骗，消灭一切偶像、魔鬼和污灵的一切权势。当我来到这里时，魔鬼因你的名从我们面前逃跑，巨龙和蛇都枯萎了，但这些人没有接受你真实的光；因此我决定按他们的愚昧降服他们。}\par

救主说：\Quote{斐理伯，因为你放弃了我这不可以恶报恶的诫命，为此你在来世将被剥夺在应许之地停留四十年；此外，这是你在此处离世离身的结局；巴尔多禄茂的份在吕考尼亚，将在那里被钉十字架；玛利亚姆将把自己的身体安放在约但河里。}\par

救主转过身来，伸开祂的手，在空中画了十字架的记号；那记号充满了光，形状如同梯子。所有下到深渊里的城中百姓，都藉着那光之十字架的梯子上来了，没有一人留在深渊中，只有暴君、司祭和他们所敬拜的毒蛇留在那里。当百姓从深渊上来时，他们抬头看见斐理伯头朝下悬吊着，巴尔多禄茂在圣殿的墙上，他们也找到了玛利亚，她恢复了起初的形状。救主在斐理伯、巴尔多禄茂、玛利亚、豹子、山羊羔、尼加诺拉和斯塔希眼前升了天；他们都以大声畏惧战兢地光荣了天主，呼喊说：\Quote{只有一位天主，祂给我们派遣了祂的救恩，这些人所宣扬的就是祂的名。因此，我们悔改我们昨日之前所陷的错误，我们不配得永生；我们相信了，因为我们亲眼看见了藉着我们而发生的奇事。}另一些人俯伏在地，朝拜宗徒们；还有人准备逃跑，说：\Quote{可能还会发生像刚才那样的地震。}\par

宗徒斐理伯头朝下悬吊着，伸出双手说：\Quote{城中百姓，请听我现在头朝下悬吊着要对你们说的话。你们已经知道天主的权能何等伟大，并看见当你们城因地震而毁灭时你们所见的奇事。这事对你们也是明显的：斯塔希的家没有被毁，他没有下到深渊，因为他相信了真天主，并接待了我们这些祂的仆人。而我，既已完成了我的天主的一切旨意，便欠祂的债，因我以恶报那对我行恶的人。}\par

那些受过洗的有些人跑去要解下头朝下悬吊着的斐理伯。他回答说：\Quote{我的弟兄们，……那些在肉体肢体上保持童贞，却在心中行奸淫，并眼目行奸淫的人，必如洪水一样泛滥。他们因听信诱人的享乐而放纵无度，忘记了认识福音的天主；他们的心充满傲慢，在敬拜中又吃又喝，忘记了圣洁的诫命，并藐视它。那世代已经偏离了正道；但那退入自己隐居处的人是有福的，因为他必在离世时得到安息。巴尔多禄茂，难道你不知道我们主的话是真生命和真知识吗？因为主在教导中对我们说：\BibleQuote{凡注视妇女，有意贪恋她的，他已经在心里奸淫了她。}\BibleRef{玛 5:28}因此，我们的弟兄伯多禄逃避一切有妇女的地方，然而他自己的女儿却成了绊脚石；他向主祈祷，主就使她的半身瘫痪，以免她受欺骗。弟兄，你看，眼目的观看带来悖逆和罪的开始，正如经上所记：\BibleQuote{女人看那棵树，见那棵树实在好看，又好吃，就被诱惑了。}\BibleRef{创 3:6}因此，童贞女的听觉应是圣洁的；她们外出时应两人同行，因为仇敌的诡计很多。她们的行走和言谈应有秩序，好能得救；否则，她们的果实就归众人所有。}\par

\Quote{我的弟兄巴尔多禄茂，把这些应许交给斯塔希，并立他为教会的治理者和主教，使他能像你一样善加教导。不要把职务托付给太年轻的人：不要立这样的人在教师的席位上，免得你亵渎基督的见证。因为教导者应当使他的行为与言语相称，好使圣言能在任何场合都自有其光荣。但我现在头朝下悬吊着，即将脱离我的身体。那么，取我的身体，用叙利亚纸预备埋葬，不要用细麻布裹它，因为人们曾用细麻布裹我们主的身体；要用纸和纸莎草紧密包裹，放在圣教会的门廊里。并为我祈祷四十天，好使天主宽恕我所犯的过犯，即我以恶报那对我行恶的人，使我在来世不必再受那四十年。}\par

说完这些话后，斐理伯祈祷说：\Quote{我的主耶稣基督，万世之父，众光之王，祢以祢的智慧使我们有智慧，祢赐给了我们高超的知识，祢仁慈地将祢良善的旨意赐给了我们，祢从未离弃我们；祢除去那些投靠祢的人的疾病；祢赐给了我们圣言，好使迷失的人归向祢；祢为小信德的人赐下征兆和奇事；祢为得胜者预备冠冕；祢是竞赛的审判者，祢赐给了我们喜乐的冠冕，祢与我们说话，使我们能抵挡那些伤害我们的人；祢是播种、收割、完成、增加并赐生命给祢所有仆人的那一位；辱骂和恐吓，因着那些藉着我们——祢的仆人——而转向祢的人，成了我们的帮助和能力。主，求祢降临，在众人面前赐给我胜利的冠冕。不要让他们的黑暗之气包围我，也不要让他们的烟火灼伤我灵魂的形像，使我得以渡过深渊的水，而不沉溺其中。我的主耶稣基督，不要让仇敌在祢——真审判者——面前找到任何可以控告我的把柄，但求祢用祢光耀的袍服披戴我，并且……}\EditorialNote{其余部分缺失。}\par

\SourceRecord{Translated by Alexander Walker.；Ante-Nicene Fathers；Vol. 8.；Edited by Alexander Roberts, James Donaldson, and A. Cleveland Coxe.；Buffalo, NY: Christian Literature Publishing Co.,；1886.；<http://www.newadvent.org/fathers/0818.htm>.}
